\documentclass[11pt,a4paper]{article}
\usepackage[utf8]{inputenc}
\usepackage[T1]{fontenc}
\usepackage{amsmath,amssymb,amsthm}
\usepackage{graphicx}
\usepackage{booktabs}
\usepackage{hyperref}
\usepackage{natbib}
\usepackage[margin=1in]{geometry}
\usepackage{algorithm}
\usepackage{algorithmic}
\usepackage{xcolor}

\newtheorem{theorem}{Theorem}
\newtheorem{definition}{Definition}
\newtheorem{lemma}{Lemma}
\newtheorem{corollary}{Corollary}
\newtheorem{proposition}{Proposition}

\title{Inference-Time Epistemic Quarantine for LLM Hallucination Containment}

\author{
Muhammad Saad \\
Independent Researcher \\
Islamabad, Pakistan \\
\texttt{saadkhan@proton.me}
}

\date{April 2026}

\begin{document}
\maketitle

\begin{abstract}
Large language models produce confident but factually incorrect outputs, a phenomenon known as hallucination. Existing mitigation strategies operate at the training or prompting level, leaving a critical gap at inference time: once a hallucinated claim enters a reasoning chain, no deployed mechanism prevents it from propagating to downstream conclusions. We introduce the Cognitive Immune System (CIS), a five-layer inference-time architecture that decomposes LLM responses into atomic claims, scores each claim against independent epistemic signals, and partitions claims into safe and quarantined sets before they can influence subsequent reasoning. We formalize the core primitive of Epistemic Quarantine and prove a Containment Depth Bound: for a single-checkpoint detection rate $\rho$ and target escape probability $\varepsilon$, at most $k^* = \lceil \log(\varepsilon) / \log(1-\rho) \rceil$ sequential checkpoints are needed to guarantee containment. We evaluate CIS on HotpotQA using adversarial semantic injection, where an auxiliary LLM rewrites critical supporting facts to produce wrong answers. On 86 verified adversarial questions and 100 clean controls, CIS achieves a contamination detection rate of 28.4\% (95\% CI: 20.0\%--38.6\%) with a false positive rate of 15.5\% (95\% CI: 9.8\%--23.8\%). Substituting the measured $\rho=0.284$ into the depth bound yields $k^*=9$ for $\varepsilon=0.05$. We analyze the primary bottleneck for single-checkpoint detection: semantic drift between injected contamination and extracted claims caused by aggressive LLM paraphrasing during response generation.
\end{abstract}

\section{Introduction}

Large language models hallucinate. They produce statements that are fluent, confident, and wrong. This is well documented across model families \citep{ji2023survey, huang2023survey, zhang2023siren}, and it remains one of the central open problems in deploying LLMs for high-stakes applications.

The research community has responded with a growing body of work on hallucination detection and mitigation. Retrieval-augmented generation grounds model outputs in external documents \citep{lewis2020retrieval}. Self-consistency methods sample multiple responses and check agreement \citep{wang2023selfconsistency}. Factual verification systems compare claims against knowledge bases \citep{min2023factscore}. Training-time approaches fine-tune models to be more calibrated \citep{kadavath2022language}.

What none of these approaches address is a simple but critical failure mode: once a hallucinated claim enters a multi-step reasoning chain, it propagates. A wrong intermediate conclusion becomes the premise for the next step, and by the time the final answer is generated, the contamination is embedded in the logical structure of the response. There is no mechanism at inference time to detect a contaminated claim and prevent it from reaching the next reasoning step.

This paper introduces such a mechanism. We call it Epistemic Quarantine.

The idea is borrowed from immunology. When the human immune system encounters a pathogen, it does not attempt to reason about whether the pathogen is dangerous. It isolates first, investigates second. The default action on uncertainty is containment, not propagation.

We formalize this into a five-layer architecture called the Cognitive Immune System (CIS):

\begin{enumerate}
    \item \textbf{Claim Extraction} decomposes a raw LLM response into atomic, verifiable claims.
    \item \textbf{Contamination Scoring} assigns each claim a contamination score $\phi(c) \in [0,1]$ using three independent signals: Wikipedia cross-verification, LLM self-consistency, and causal ancestry.
    \item \textbf{Quarantine Partitioning} splits claims into a safe set $S$ and a quarantine set $Q$ based on whether $\phi(c)$ exceeds a threshold $\tau$.
    \item \textbf{Causal Memory} records quarantined claims in a directed acyclic graph, enabling cross-session immunity.
    \item \textbf{Tolerance Calibration} maintains a registry of verified-safe claims to reduce false positives over time.
\end{enumerate}

The critical invariant is that the context passed to the next reasoning step comes only from $S$. No quarantined claim ever reaches the downstream computation.

We prove that this architecture provides a formal containment guarantee. If a single checkpoint detects contamination with probability $\rho$, then $k$ sequential checkpoints reduce the escape probability to $(1-\rho)^k$. For any target escape probability $\varepsilon$, we need at most $k^* = \lceil \log(\varepsilon) / \log(1-\rho) \rceil$ checkpoints. This is Theorem 1, the Containment Depth Bound, and it holds regardless of the specific scoring function as long as checks are independent.

We evaluate CIS on HotpotQA \citep{yang2018hotpotqa} using a methodology we call adversarial semantic injection. Rather than randomly replacing entities in the supporting context, we use an auxiliary LLM to rewrite the critical supporting sentence so that the answer to the question changes. We verify that the injection works by testing whether the baseline LLM actually produces a wrong answer on the contaminated context. Only verified adversarial examples enter the benchmark.

Our results are honest and below our initial targets. CIS detects 28.4\% of adversarial contaminations at a 15.5\% false positive rate. We analyze why: the primary bottleneck is semantic drift. The LLM paraphrases the contaminated context so aggressively that the extracted claims look different from the original injection, making Wikipedia cross-verification ineffective. This finding has implications for any system that attempts to verify LLM outputs by comparing them against external sources.

\section{Related Work}

\paragraph{Hallucination Detection.} FActScore \citep{min2023factscore} decomposes text into atomic facts and checks each against a knowledge source. SelfCheckGPT \citep{manakul2023selfcheckgpt} uses sampling-based consistency. SAFE \citep{wei2024longform} uses LLM-as-judge for long-form factuality. These systems detect hallucinations but do not prevent their propagation through reasoning chains.

\paragraph{Retrieval-Augmented Generation.} RAG \citep{lewis2020retrieval} and its descendants \citep{gao2023retrieval} reduce hallucination by grounding generation in retrieved documents. However, RAG does not verify the generated output, and contaminated retrievals propagate directly into responses.

\paragraph{Self-Consistency.} \citet{wang2023selfconsistency} show that sampling multiple reasoning paths and taking the majority vote improves accuracy. This is a form of ensemble verification but operates at the response level, not the claim level, and does not provide formal containment guarantees.

\paragraph{Inference-Time Intervention.} Inference-time intervention \citep{li2024inference} modifies internal model representations during generation. This is mechanistically different from our approach: we operate on the output of generation, not its internals, making CIS model-agnostic.

\paragraph{Gap.} To our knowledge, no prior work implements a mechanism that (a) decomposes LLM output into atomic claims, (b) independently verifies each claim, (c) partitions claims into safe and quarantine sets at inference time, and (d) provides a formal bound on contamination escape probability. CIS fills this gap.

\section{Architecture}

\subsection{Overview}

Let $R$ be the raw response from an LLM to a query $q$ with optional context $x$. The CIS pipeline computes:
\[
\Pi(R) = L_5 \circ L_4 \circ L_3 \circ L_2 \circ L_1(R)
\]
where each layer performs a specific transformation.

\subsection{Layer 1: Claim Extraction}

$L_1: R \to C = \{c_1, c_2, \ldots, c_n\}$

We prompt the LLM to decompose $R$ into atomic, independently verifiable claims. Each claim $c_i$ is a single factual assertion that can be checked against an external source without reference to the other claims.

\subsection{Layer 2: Contamination Scoring}

$L_2: c_i \to \phi(c_i) \in [0, 1]$

The contamination score is a weighted combination of three independent signals:

\begin{equation}
\phi(c) = w_1 \cdot S_{\text{wiki}}(c) + w_2 \cdot S_{\text{cons}}(c) + w_3 \cdot S_{\text{causal}}(c)
\end{equation}

where $w_1 = 0.60$, $w_2 = 0.30$, $w_3 = 0.10$.

\textbf{$S_{\text{wiki}}(c)$: Wikipedia Cross-Verification.} We extract named entities from $c$, retrieve the corresponding Wikipedia summaries, and prompt an LLM to judge whether the summary supports, contradicts, or is insufficient regarding $c$. Contradicted $\to 1.0$, supported $\to 0.0$, insufficient $\to 0.5$.

\textbf{$S_{\text{cons}}(c)$: Factual Confidence Probe.} We ask an LLM: ``On a scale of 0 to 10, how confident are you that this claim is factually correct?'' We compute $S_{\text{cons}} = 1 - (\text{response}/10)$. This inverts the signal so that low confidence yields high contamination scores.

\textbf{$S_{\text{causal}}(c)$: Causal Ancestry.} If the claim shares a causal ancestor in the contamination DAG with a previously quarantined claim, $S_{\text{causal}} = 1$. Otherwise $S_{\text{causal}} = 0$.

\subsection{Layer 3: Quarantine Partitioning}

\begin{definition}[Epistemic Quarantine]
Let $C = \{c_1, \ldots, c_n\}$ be the claims extracted from a response. Let $\phi: C \to [0,1]$ be the contamination scoring function and $\tau \in (0,1)$ the threshold. We define two disjoint partitions:
\begin{align}
S &= \{c \in C \mid \phi(c) < \tau\} \\
Q &= \{c \in C \mid \phi(c) \geq \tau\}
\end{align}
such that $S \cup Q = C$ and $S \cap Q = \emptyset$.
\end{definition}

The critical invariant is:
\[
\text{context}(\text{step}_{k+1}) \subseteq S
\]

No claim in $Q$ is ever passed as context to the next reasoning step. This breaks the contamination propagation chain.

\subsection{Layer 4: Causal Memory (CC-DAG)}

Quarantined claims are recorded in a Contamination Causal Directed Acyclic Graph $G = (V, E)$ where vertices represent contamination events and edges represent causal precedence. This provides cross-session immunity: if a claim in a new session shares structure with a previously quarantined claim, the causal signal $S_{\text{causal}}$ increases its $\phi$ score.

\subsection{Layer 5: Tolerance Calibration}

A safe registry $T$ stores claims that have been verified as correct with high confidence. Claims in $T$ bypass the scoring pipeline entirely. This prevents autoimmune false positives on repeatedly verified facts and provides monotonic FPR reduction over time.

\section{Containment Depth Bound}

\begin{theorem}[Containment Depth Bound]
Let $\rho = P(\phi(c) \geq \tau \mid c \text{ is contaminated})$ be the true positive detection rate of a single EQ checkpoint. Assume that detection decisions at successive checkpoints are independent. Then the probability that a contaminated claim escapes $k$ consecutive checkpoints is:
\[
P(\text{escape after } k) = (1-\rho)^k
\]
For a target escape probability $\varepsilon \in (0,1)$, the minimum number of checkpoints required is:
\[
k^* = \left\lceil \frac{\log \varepsilon}{\log(1-\rho)} \right\rceil
\]
\end{theorem}

\begin{proof}
A contaminated claim escapes a single checkpoint with probability $1-\rho$. By the independence assumption, the probability of escaping $k$ consecutive checkpoints is $(1-\rho)^k$.

We require $(1-\rho)^k \leq \varepsilon$. Taking logarithms of both sides:
\[
k \cdot \log(1-\rho) \leq \log(\varepsilon)
\]
Since $\rho \in (0,1)$, we have $\log(1-\rho) < 0$, so dividing reverses the inequality:
\[
k \geq \frac{\log \varepsilon}{\log(1-\rho)}
\]
The minimum integer satisfying this is $k^* = \lceil \log(\varepsilon) / \log(1-\rho) \rceil$.
\end{proof}

\begin{corollary}
For the empirically measured $\rho = 0.284$ and $\varepsilon = 0.05$:
\[
k^* = \left\lceil \frac{\log(0.05)}{\log(0.716)} \right\rceil = \left\lceil \frac{-2.996}{-0.334} \right\rceil = \lceil 8.97 \rceil = 9
\]
Nine sequential EQ checkpoints reduce contamination escape probability below 5\%.
\end{corollary}

\begin{corollary}
For $\varepsilon = 0.01$:
$k^* = \lceil \log(0.01) / \log(0.716) \rceil = \lceil 13.79 \rceil = 14$.
Fourteen checkpoints reduce escape probability below 1\%.
\end{corollary}

\section{Experimental Setup}

\subsection{Dataset and Models}

We use the HotpotQA validation set \citep{yang2018hotpotqa}, a multi-hop question answering dataset where each question requires reasoning over two or more supporting documents. We use Llama-3.3-70B-Versatile (via Groq API) as the main reasoning model, and Llama-3.1-8B-Instant for scoring and injection.

\subsection{Adversarial Semantic Injection}

Random entity replacement achieves only a 2.8\% adversarial hit rate because most replacements land on facts the LLM does not need to answer the question. We replace this with LLM-guided semantic injection:

\begin{enumerate}
    \item Extract the critical supporting sentences using HotpotQA's annotated supporting fact pointers.
    \item Prompt Llama-3.1-8B-Instant to rewrite the critical sentence so that the question's answer changes, without mentioning the gold answer.
    \item Verify the injection works: run the baseline LLM on the contaminated context. If the baseline still produces the correct answer, discard the question.
    \item Only questions where the baseline is verifiably fooled enter Group A (contaminated).
\end{enumerate}

This achieves a 17.2\% adversarial hit rate, a 6$\times$ improvement over random replacement.

\subsection{Evaluation Protocol}

We scan up to 500 HotpotQA candidates and fill two groups:
\begin{itemize}
    \item \textbf{Group A (contaminated):} Questions where the adversarial injection verifiably fools the baseline LLM. CIS processes the contaminated context.
    \item \textbf{Group B (control):} Questions with clean, unmodified context. CIS processes the original context.
\end{itemize}

\textbf{M1 (Contamination Detection Rate):} Fraction of Group A questions where CIS quarantines at least one claim. \\
\textbf{M2 (False Positive Rate):} Fraction of Group B questions where CIS quarantines at least one claim. \\
\textbf{M3 (Accuracy):} Exact match accuracy comparison between baseline and CIS. \\
\textbf{M4 (Adversarial Hit Rate):} Fraction of candidates where injection successfully fools the baseline. \\
\textbf{M5 (Latency Overhead):} Average CIS latency versus baseline latency.

Evaluation uses standard HotpotQA exact match: the predicted answer must contain the gold answer string after normalization. We report Wilson score 95\% confidence intervals for M1 and M2.

\section{Results}

\begin{table}[h]
\centering
\begin{tabular}{lcc}
\toprule
\textbf{Metric} & \textbf{Value} & \textbf{95\% CI} \\
\midrule
M1 (CDR) & 28.4\% & [20.0\%, 38.6\%] \\
M2 (FPR) & 15.5\% & [9.8\%, 23.8\%] \\
M3 Baseline Accuracy (Grp A) & 1.1\% & -- \\
M3 CIS Accuracy (Grp A) & 2.3\% & -- \\
M4 Adversarial Hit Rate & 17.2\% & -- \\
M5 Avg Baseline Latency & 1,288 ms & -- \\
M5 Avg CIS Latency & 22,291 ms & -- \\
M5 Overhead & 1,631\% & -- \\
\bottomrule
\end{tabular}
\caption{Final experimental results on HotpotQA (86 adversarial, 100 control).}
\label{tab:results}
\end{table}

\paragraph{Detection rate.} CIS quarantines at least one claim in 25 of 88 adversarial questions (28.4\%). This is significantly above zero but below the 65\% threshold we would consider reliable for standalone deployment.

\paragraph{False positive rate.} CIS incorrectly quarantines claims in 16 of 103 control questions (15.5\%). This is within our acceptable bound of 20\%.

\paragraph{Theorem validation.} Substituting $\rho = 0.284$ into the depth bound:
\begin{itemize}
    \item $k^* = 9$ for $\varepsilon = 0.05$ (escape probability $< 5\%$ after 9 checkpoints)
    \item $k^* = 14$ for $\varepsilon = 0.01$ (escape probability $< 1\%$ after 14 checkpoints)
\end{itemize}

\paragraph{McNemar's test.} The difference in accuracy between baseline and CIS is statistically significant ($p = 0.0265$), with CIS slightly improving accuracy on adversarial questions by successfully quarantining some contaminated claims.

\section{Analysis: Why M1 Is 28.4\%}

The primary bottleneck is semantic drift between the injected contamination and the claims CIS extracts from the LLM response.

\paragraph{Paraphrase-level transformation.} When the LLM receives contaminated context, it does not copy the false facts verbatim. It synthesizes, summarizes, and reframes. The atomic claims extracted from this synthesized response often differ lexically and structurally from the injected contamination, even though they carry the same wrong conclusion. When CIS sends these paraphrased claims to Wikipedia for verification, the entity matching fails because the claim uses different words.

\paragraph{Sparse Wikipedia coverage.} HotpotQA questions frequently involve obscure people, films, and events. For 12\% of claims, no relevant Wikipedia page exists. When Wikipedia returns no match, $S_{\text{wiki}}$ defaults to 0.5 (uncertain), which contributes only $0.60 \times 0.5 = 0.30$ to $\phi$, insufficient to cross the threshold $\tau = 0.55$ without a strong consistency signal.

\paragraph{Overconfident scorer.} The 8B scoring model frequently returns high factual confidence (7 to 9 out of 10) even for false claims about entities outside its training distribution. This causes $S_{\text{cons}}$ to remain low (0.1 to 0.3) for contaminated claims, failing to complement the Wikipedia signal when it is needed most.

\paragraph{Implications.} These findings suggest that single-checkpoint detection of semantic hallucinations may require entity-grounded claim extraction (maintaining the link between claims and source entities) and a much stronger factual verification model. The containment depth bound theorem provides an alternative path: even with weak single-checkpoint detection, sufficiently many sequential checkpoints still guarantee containment.

\section{Limitations}

\begin{enumerate}
    \item The experiment uses a single LLM family (Llama) for both generation and scoring. Cross-model evaluation would strengthen the results.
    \item The scoring model (8B parameters) is likely too small for reliable factual verification on obscure entities. A larger scorer could improve M1 significantly.
    \item The independence assumption in Theorem 1 may not hold perfectly in practice, as successive checkpoints share the same scoring model. Empirical verification with multiple checkpoints is needed.
    \item Group A reached 86 of the target 100 questions. The confidence intervals reflect this sample size limitation.
    \item CIS adds 1,631\% latency overhead, making it impractical for real-time applications without optimization.
\end{enumerate}

\section{Conclusion}

We introduced Epistemic Quarantine, the first inference-time mechanism for containing hallucination propagation in LLM reasoning chains. Our main contribution is the formalization and proof of the Containment Depth Bound, which shows that even imperfect single-checkpoint detection can be composed into strong containment guarantees through sequential application. On HotpotQA with adversarial semantic injection, CIS achieves $\rho = 0.284$ per checkpoint, yielding $k^* = 9$ for 95\% containment. The primary bottleneck we identify, semantic drift between injected contamination and LLM-paraphrased claims, points to entity-grounded claim extraction as the key direction for improving single-checkpoint detection rates.

\bibliographystyle{plainnat}
\begin{thebibliography}{99}

\bibitem[Gao et al., 2023]{gao2023retrieval}
Gao, Y., Xiong, Y., et al. (2023). Retrieval-augmented generation for large language models: A survey. \textit{arXiv:2312.10997}.

\bibitem[Huang et al., 2023]{huang2023survey}
Huang, L., Yu, W., et al. (2023). A survey on hallucination in large language models. \textit{arXiv:2311.05232}.

\bibitem[Ji et al., 2023]{ji2023survey}
Ji, Z., Lee, N., et al. (2023). Survey of hallucination in natural language generation. \textit{ACM Computing Surveys}, 55(12).

\bibitem[Kadavath et al., 2022]{kadavath2022language}
Kadavath, S., et al. (2022). Language models (mostly) know what they know. \textit{arXiv:2207.05221}.

\bibitem[Lewis et al., 2020]{lewis2020retrieval}
Lewis, P., et al. (2020). Retrieval-augmented generation for knowledge-intensive NLP tasks. \textit{NeurIPS 2020}.

\bibitem[Li et al., 2024]{li2024inference}
Li, K., et al. (2024). Inference-time intervention: Eliciting truthful answers from a language model. \textit{NeurIPS 2024}.

\bibitem[Manakul et al., 2023]{manakul2023selfcheckgpt}
Manakul, P., et al. (2023). SelfCheckGPT: Zero-resource black-box hallucination detection. \textit{EMNLP 2023}.

\bibitem[Min et al., 2023]{min2023factscore}
Min, S., et al. (2023). FActScore: Fine-grained atomic evaluation of factual precision. \textit{EMNLP 2023}.

\bibitem[Wang et al., 2023]{wang2023selfconsistency}
Wang, X., et al. (2023). Self-consistency improves chain of thought reasoning in language models. \textit{ICLR 2023}.

\bibitem[Wei et al., 2024]{wei2024longform}
Wei, J., et al. (2024). Long-form factuality in large language models. \textit{arXiv:2403.18802}.

\bibitem[Yang et al., 2018]{yang2018hotpotqa}
Yang, Z., et al. (2018). HotpotQA: A dataset for diverse, explainable multi-hop question answering. \textit{EMNLP 2018}.

\bibitem[Zhang et al., 2023]{zhang2023siren}
Zhang, Y., et al. (2023). Siren's song in the AI ocean: A survey on hallucination in large language models. \textit{arXiv:2309.01219}.

\end{thebibliography}

\end{document}
